\chapter{Methodology}
\label{ch:models}

% razionale (obiettivo, criteri, regole di inclusione) + famiglie di approcci con pro e contro
This chapter sets out the rationale of the study (\cref{sec:meth:rationale}) and the
approaches that were compared, in five families: causal rules on the price, stateless
learners on hand-made features, sequence models of the order-book stream, graph models,
and sources of information outside the book.  \todo{each family closes with its
advantages and limits a priori; a summary table at the end}

\section{Objective, criteria and rules of inclusion}
\label{sec:meth:rationale}

\subsection{Objective}
\label{sec:meth:objective}

The object of the study is the direction of the mid price of a Polymarket market over a
horizon $H$, read as the three-class label of \cref{sec:data:labels}, at every block in
which the market updated its book, for five horizons from one minute to twelve hours.
The task is the prediction of the price while the market is open, not of the outcome of
the event; the measure is balanced accuracy, because the three classes are of very
unequal size and change size with the horizon; and the unit of every comparison is the
paired difference between two models on the same samples, with an interval.

\subsection{Criteria of success}
\label{sec:meth:criteria}

A model is said to read information that another does not when its balanced accuracy
on the validation day exceeds the other's on the same samples, after the same post-hoc
calibration, by a margin whose 95\% hour-block bootstrap interval excludes zero.  The
reference competitor is first the best causal rule found on the data and then, once it
exists, the best stateless learner on hand-made features.  \todo{ties; the two views of
the samples (full stream and deduplicated); the locked test period}

\subsection{Rules of inclusion and exclusion}
\label{sec:meth:rules}

\todo{to expand, one paragraph each:}
\begin{enumerate}[nosep]
  \item \textbf{Causal inputs only.}  Everything an approach reads at time $t$ is
  computed from data strictly before $t$, and every label from data after it.
  \item \textbf{Public information that the dataset contains}, plus two public external
  sources that could be aligned to it (the news flow, a pretrained forecaster); no
  external reference prices.
  \item \textbf{One model for all markets}, never one model per market.
  \item \textbf{Every approach is measured against the best simpler one on the same
  inputs}, and where possible contains it as a special case.
  \item \textbf{Families that span the questions}: each research question has a paired
  comparison.  Out of scope: outcome forecasting with language models, trading
  strategies, per-market models.
  \item \textbf{A fixed compute budget}: one run per GPU, hours per run; short tests to
  screen the variants, full-length runs behind every claim.
\end{enumerate}

\section{Baselines: causal rules on the price}
\label{sec:models:baselines}

All baselines are causal rules on the mid price of the market itself.  With
$\Delta^{-}_{L} = \operatorname{mid}(t) - \operatorname{mid}(t-L)$ the change over the
past lookback $L$:
\begin{itemize}[nosep]
  \item \textbf{Constant}: always \textsc{const}; balanced accuracy $1/3$ by construction.
  \item \textbf{Momentum-$L$}: predict $\operatorname{sign}(\Delta^{-}_{L})$.
  \item \textbf{Mean Reversion-$L$}: predict $-\operatorname{sign}(\Delta^{-}_{L})$, and
  \textsc{const} if the price did not move.
\end{itemize}

\paragraph{Advantages and limits.}  \todo{to write}

\section{Stateless learners on hand-made features}
\label{sec:models:stateless}

A LightGBM classifier \cite{ke2017lightgbm} on the sixty-one features of
\cref{sec:prep:features}, trained with class-balanced weights on 7--11 March and
early-stopped on 12 March, adding one group of features at a time.  On the same table,
an MLP and a TabM ensemble, trained as tabular models should be (\cref{sec:bg:tabular}),
to tell whether a difference between the trees and a deep model is the learner or the
training regime.  \todo{hyper-parameters}

\paragraph{Advantages and limits.}  \todo{to write}

\section{Sequence models of the order-book stream}
\label{sec:models:sequence}

\subsection{The continuous-time stateful model}
\label{sec:models:ctdgnn}

The model (CT-DGNN in the code) processes the stream block by block.  At block time $t$
it receives the markets that updated their book since the previous block, with a window
of their last $K$ book states, and the transactions of the block as edges between
wallets and markets; it keeps a state vector for every market and every wallet.
\todo{one step of the model, with the architecture figure from the old chapter:}
\begin{itemize}[nosep]
  \item states and initialisation; continuous-time decay of the states;
  \item order-book update: the window encoders (space first, or time first);
  \item hierarchical memory at three rates; global state;
  \item long-lookback price features and the direct path, which contains the
  mean-reversion rule as a special case;
  \item predictor, loss and training (truncated BPTT, one pass through the stream).
\end{itemize}

\subsection{The stateful model as the residual of the trees}
\label{sec:models:residual}

\todo{cross-fitted tree log-probabilities added to the output logits; at initialisation
the model is the tree model, and the number that measures the state is the gain over it}

\subsection{A pretrained forecaster used zero-shot}
\label{sec:models:zeroshot}

\todo{Chronos-2 on the last 512 mids of the market; three decision rules; the same
quantiles as features of the stateful model}

\paragraph{Advantages and limits.}  \todo{to write}

\section{Graph models}
\label{sec:models:graph}

\subsection{Message passing on transactions}
\label{sec:models:mp}

\todo{edges $(w, m)$ with a 7-dimensional feature; attention over the incident edges
and a GRU update; the token-orientation feature; the enrichments (wallet skill, placebo
skill, past behaviour); the ablation without edges}

\subsection{The two graphs}
\label{sec:models:twographs}

\todo{the wallet--market graph, read by message passing, and the market--market graph
between the contracts of the same event, measured first through the sibling features}

\paragraph{Advantages and limits.}  \todo{to write}

\section{Information outside the book}
\label{sec:models:external}

\todo{the foundation-model forecasts, the news flow, the text of the market; each
measured first in the stateless learner}

\paragraph{Advantages and limits.}  \todo{to write}

\section{The families at a glance}
\label{sec:models:summary}

\todo{table: family, members, what it reads, where its state lives, reference}
